\subsection{滑らかな楕円表示での有限次数開条件}
\label{subsec:cubic-CY-three-alpha-elliptic-relative}

\begin{theorem}\label{thm:cubic-CY-three-alpha-elliptic-relative-smooth-open}
滑らかな先頭ポテンシャルと滑らかな楕円境界を持つHesse表示の
全低次係数および有限中心値の係数空間を取る。
この空間には、有限行列のrankで定まる開集合 $\mathcal W$ があり、
次の条件を満たす。
\begin{enumerate}
\item $\mathcal W$ 上のアフィン切断は相対的に双加群滑らかであり、
対角は長さ2以下の有限生成射影解消を持つ。
これは任意底変更で保たれ、幾何学的fiberは全てH3である。
\item 境界の平行移動が無限位数の幾何学的点では、
$\mathcal W$ に入ることと双加群滑らかさは同値である。
従ってこの範囲の双加群滑らかな点には、
有限位数の方向も許す代数的なH3表示近傍が存在する。
\end{enumerate}
\end{theorem}
\begin{proof}
$p=1$ に規格化し、滑らかさを満たす先頭二係数の開集合に、
対称な六つの一次係数、三つの定数係数、切断値を加える。
全包絡は補題\ref{lem:sklyanin-central-PBW-CY}のPBW表示であり、
各次数の関係挿入行列はこの被約係数空間で一定階数を持つ。
従って斉次成分は有限局所自由で、任意底変更と可換する。
定理\ref{thm:cubic-CY-three-alpha-smooth-boundary}により中央三次の
核も一定階数2を持つので、底を局所化して先頭 $g$ の持ち上げ $\omega$ を選ぶ。
商の各次数も、全幾何fiberで三次正則元を割るため有限局所自由である。
従ってアフィン包絡 $D$ と切断 $U_\lambda$ の次数フィルターも
有限局所自由な各項と各随伴成分を持つ。

基礎係数環を $R$ と書く。$D\otimes_R D^{\mathrm{op}}$ 上では、
絶対対角の小解消を生成元の差、三関係の非可換微分、
生成元の差から作る。最後の合成が零であることは巡回微分恒等式であり、
随伴複体は斉次Koszul双加群解消である。
固定したフィルター次数の複体は有限局所自由で、幾何学的fiber上で完全である。
従って局所的に分裂完全で、差
$\omega_L-\omega_R$ の高次ホモトピーを重み $3a$ の範囲内で解ける。
必要な未知数は、長さ3の解消で添字 $a$ が1または2のものに限られる。
中央差によるShamash解消を取って $R[\omega]$ 上の相対対角解消を作り、
さらに $\omega=\lambda$ で底変更した解消を $K_\lambda$ と書く。
$D$ の $R[\omega]$ 上の平坦性は、先頭の $R[g]$ 上で
次数別に $gS$ の補空間を取って持ち上げることから従う。
従って $K_\lambda$ は $U_\lambda$ の対角を解消し、
この構成の完全性は任意底変更で保たれる。

相対対角のコチェインを $C$、切断対角のコチェインを
$C_\lambda=\operatorname{Hom}_{U_\lambda^e}(K_\lambda,U_\lambda)$ と書く。
内部重み高々0の部分複体 $F_0C_\lambda$ の各項は有限局所自由である。
この複体のコホモロジーは解消の局所選択に依存しない。
重み付き自由基底を用い、strictな完全性で重みを保つ比較写像を帰納的に持ち上げると、
逆比較と二つのホモトピーも重みを保って取れる。
従って同じフィルター部分で連鎖ホモトピー同値が成り立つ。
四次コホモロジーが幾何学的fiberで零であることは
\[
 \operatorname{rank} d^3+\operatorname{rank}d^4=\operatorname{rank}(F_0 C_\lambda^4)
\]
であり、二つの階数の下界で与える開条件である。
これを $\mathcal W$ とする。$\mathcal W$ 上では
行列の単元minorを取って、まず $d^3$ の像を部分束として分離し、
残る直和因子を $d^4$ で分裂単射にする。
従って四次の完全性は任意の底変更で保たれる。

変換射の二乗は、Shamash添字を二つ下げる鎖射と増大写像の合成で
表されるため、内部重み高々 $-6$ である。
従って $\mathcal W$ 上でこの類は境界となる。
中央差による二段三角を用いると、対角は完全複体の直和因子となり、
相対双加群滑らかさを得る。
相対絶対小解消の自己双対性は巡回ポテンシャルの左・右微分から従う。
二つの中心正則式 $\omega_L-\lambda,\omega_R-\lambda$ で双対移行すると、
切断の対角の双対は次数2に集中する。
有限射影解消の最後の双対微分を分裂させる議論で、解消は長さ2以下に短縮できる。
加群と各射影項は底上平坦なので、任意底変更で完全性と射影性を保つ。
幾何学的fiberのH3性は
定理\ref{thm:cubic-CY-three-alpha-smooth-boundary}のアフィン認識による。

最後に無限位数の幾何学的点で検査する。中央切断の先頭は次数3なので
コチェインには短完全列
\[
 0\longrightarrow F_{-3}C
   \xrightarrow{\omega-\lambda} F_0C
   \longrightarrow F_0 C_\lambda\longrightarrow0
\]
がある。定理\ref{thm:cubic-CY-three-alpha-elliptic-diagonal-test}の
最高重みの降下により、
$H^4(F_{-3}C)=H^4(F_0C)=V$ かつ $H^5(F_{-3}C)=0$ である。
よって $H^4(F_0C_\lambda)=V/(\omega-\lambda)V$ となり、
同定理の必要十分条件と一致する。
\end{proof}

\begin{remark}
有限位数の全H3点や、無限位数で判別式の根となる点に対する
開性はこの有限次数判定の結論に含めない。
\end{remark}

\begin{corollary}\label{cor:cubic-CY-three-alpha-elliptic-strong-noetherian}
このHesse表示の斉次切断環は、任意のnoetherian底変更上で左右noetherianである。
特に $\mathcal W$ の表示から得る構造対象付きtails圏には、
命題\ref{prop:h3-presentation-canonical-degree-descent}と
\ref{prop:h3-presented-qgr-fpqc}の有効降下が適用できる。
\end{corollary}
\begin{proof}
斉次切断環を $R_{\rm gr}$ とし、中央一次元 $t$ による商を取る。
これは相対楕円曲線上の捩れ斉次座標環である。
実際、滑らかな楕円上で点関係行列の階数は2であり、
左右の核から相互逆の点対応を得る。
この構成は有限ランクの核であり、普遍係数底上で形成した後の任意底変更に両立する。
したがって平行移動と次数3の線型束から相対捩れ環を作る。
正次数の線型束の $H^1$ は各fiberで零であり、切断形成が底変更に両立する。
標準生成元による元の商環からの写像は各幾何学的fiberで同型で、
両辺は各次数局所自由なので、相対的にも同型である。

任意のnoetherian底変更上で、捩れ積の線型束 $\mathcal L_n$ のfiber次数は $3n$ である。
任意の固定線型束 $\mathcal H$ を掛けても $n$ を大きくすれば
全fiberで次数が正になり、相対的にampleである。
よって \cite[定理1.3, 7.14]{KeelerAmpleFilters} により、捩れ環は左右noetherian。
$t$ は次数正の中央正則元なので、$t$ に関する随伴環は
$(R_{\rm gr}/(t))[T]$ であり、次数別に $t$ 進分離的である。
斉次イデアルの最高 $t$ 可除次数を順に上げる標準的な帰納法で有限生成性を持ち上げ、
次に通常の次数フィルターで任意のイデアルにも有限生成性を得る。
この議論は全て底変更後にも成り立つ。
\end{proof}

\begin{proposition}\label{prop:cubic-CY-three-alpha-relative-section-arrows}
このHesse表示の任意のnoetherian底変更上で、
斉次切断環は $(\qgr R_{\rm gr},\pi R_{\rm gr},(1))$ のsection環として回復される。
shiftに可換する構造対象付きexact同値の全矢は、
切断環の次数付き同型を共通次数スカラーで割ることで得られる。
ここでshiftとの可換同型は存在を要求し、矢の追加標識にはしない。
\end{proposition}
\begin{proof}
普遍係数底では、二次包絡の双対関係行列の次数別rankが各点で一定である。
係数底が被約なので、Koszul双対の次数別成分とその双対 $W_i$ は
有限局所自由であり、任意底変更と両立する。
底のaugmentation加群には、$W_i$ を使った相対Koszul解消がある。
各内部次数で幾何学的fiber上完全な有限局所自由複体なので、
解消と、包絡を係数としたその双対のGorenstein消滅は
非被約底上も含めた任意底変更で成立する。
正則三次を割る次数1の変換により、
$\underline{\operatorname{Ext}}^j_{R_{\rm gr}}(\mathcal O_T,R_{\rm gr})=0$
が $j=0,1,2$ で成り立つ。
$R_{\rm gr}/(R_{\rm gr})_{\ge n}$ は次数別の有限局所自由augmentation加群を
商に持つ有限濾過を持つ。
従ってこれを第一引数とした同じ低次Extも零となる。
$0\to (R_{\rm gr})_{\ge n}\to R_{\rm gr}\to R_{\rm gr}/(R_{\rm gr})_{\ge n}\to0$
で $\operatorname{Hom}(-,R_{\rm gr}(\ell))$ を計算し、飽和sectionの直極限を取ると、
全整数 $\ell$ でsectionの自然写像は同型である。

exact同値 $F$ とshiftとの可換同型を一つ取れば、
section積を保つ次数付き同型が得られる。
これは全ての $\pi R_{\rm gr}(i)$ 間の射を再現し、
これらの有限直和による表示を取ると元のexact同値を自然同型の範囲で再現する。
逆に次数付き同型からtails同値が得られる。
恒等関手の自然自己同型は各 $\pi R_{\rm gr}(i)$ 上で底の単元である。
隣接する対象間の射は有限局所自由で正の階数を持つので、
自然性から隣接する単元は等しい。有限表示によって全成分で同じ単元となる。
従ってshift比較の選択は共通の次数スカラーを変える。
次数付き同型がtails上で恒等と自然同型になる場合も、
隣接する対象の乗法を比較すれば一次上でただ一つのスカラー倍となり、
一次生成性で全次数のスカラー倍に限られる。
逆に次数スカラーは各次数への同じ倍率の累乗で恒等と自然同型になる。
従って主張の矢が得られる。二次包絡と切断直線はそれぞれ二次関係と
三次核から関手的に回復される。
\end{proof}
